\section{PWA and Monitoring System}
\label{sec:ch3-pwa-monitoring}

The monitoring system was implemented as a PWA connected to Firebase. The PWA
was not placed in the protection path; the ESP32-S3 still made relay decisions
locally. The PWA instead provided live measurements, protection status, event
history, CSV download, labeling support, and review of waveform-derived
features. This separation allowed the prototype to remain protective during
network loss while still supporting data collection and analysis.

Figure~\ref{fig:ch3-pwa-block} shows the PWA data path.
\begin{figure}[H]
  \centering
  \includegraphics[width=0.82\textwidth]{ch3/figure_3.12.png}
  \caption{PWA block diagram and Firebase data path.}
  \label{fig:ch3-pwa-block}
\end{figure}

Figure~\ref{fig:ch3-pwa-block} shows that the firmware publishes selected
telemetry to Firebase, while the browser-based PWA subscribes to the same data
for display and logging. This arrangement matches the PWA role described in
Firebase documentation: the web client improves access and installation
behavior, but the embedded device remains responsible for field decisions
[55], [56].

The live dashboard used during testing is shown in
Figure~\ref{fig:ch3-pwa-dashboard}.
\begin{figure}[H]
  \centering
  \includegraphics[width=0.92\textwidth]{ch3/figure_3.13.png}
  \caption{Actual PWA live dashboard.}
  \label{fig:ch3-pwa-dashboard}
\end{figure}

Figure~\ref{fig:ch3-pwa-dashboard} displays voltage, current, direct NTC
temperature, estimated socket temperature, THD, spectral features, protection
state, and relay state. The dashboard intentionally distinguishes direct NTC
temperature from estimated socket temperature because the relay-facing thermal
decision uses the compensated estimate, not the raw thermistor reading alone.

The event-history page is shown in Figure~\ref{fig:ch3-pwa-history}.
\begin{figure}[H]
  \centering
  \includegraphics[width=0.92\textwidth]{ch3/figure_3.14.png}
  \caption{Actual PWA history-log view.}
  \label{fig:ch3-pwa-history}
\end{figure}

Figure~\ref{fig:ch3-pwa-history} documents the stored-event interface used to
review warnings, trips, and data records after tests. The history function was
also useful for checking whether the firmware sequence matched the expected
warning, persistence, and latching behavior during overload, heating, and
arc-fault trials.

\section{Telemetry Fields}
\label{sec:ch3-telemetry-fields}

The PWA displayed only a selected subset of firmware variables. The most
important telemetry fields are listed in Table~\ref{tab:ch3-telemetry}.
\begin{table}[H]
  \centering
  \caption{Selected telemetry fields shown in the PWA}
  \label{tab:ch3-telemetry}
  \begin{tabular}{p{0.30\textwidth}p{0.25\textwidth}p{0.35\textwidth}}
    \toprule
    Field & Firmware source & Interpretation \\
    \midrule
    Voltage and current & Calibrated sensor routines & Electrical state of the connected load \\
    NTC temperature & Thermistor routine & External temperature at the sensor position \\
    Estimated socket temperature & Thermal estimator & Protection-facing estimate of socket temperature \\
    Arc probability and class & Random Forest inference & TinyML arc-fault decision support \\
    Context family and confidence & Context prototype model & Load-family estimate used by arc-threshold selection \\
    Relay and fault state & Protection manager & Whether the load is energized or latched off \\
    \bottomrule
  \end{tabular}
\end{table}

Table~\ref{tab:ch3-telemetry} shows that the PWA was both an operator display
and a research instrument. During data collection, these fields allowed the
researchers to compare raw measurements, computed estimates, and protection
outputs in the same session record.
